% source: manti.01.ma, 1ère SM, série fonctions (Drive 0B6pfSJcEVng0b3JzWGF4bnZPZk0).
% Ex.2 = dGtB Ex.2 ; Ex.3, 5, 6 = fonctions_s1 ; Ex.8, 9, 11 = fonctions_s2 -> non repris
\ex{Fonctions}{b3Jz}{Ex. 1}{6}
On considère la fonction $f$ définie par $f(x)=\dfrac{2x^2+7x+7}{x^2+3x+3}$.
\begin{enumerate}
\item Déterminer l'ensemble de définition $D_f$ de la fonction $f$.
\item Montrer que $\max f(x)=\dfrac73$ et $\min f(x)=1$.
\end{enumerate}
\fin
\ex{Fonctions}{b3Jz}{Ex. 4}{6}
Soit $f$ la fonction définie sur $\mathbb{R}$ par $f(x)=\dfrac{3x-4}{\sqrt{x^2+1}}$.
\begin{enumerate}
\item Montrer que $(\forall x\in\mathbb{R})\ |f(x)|\leq5$.
\item La fonction admet-elle une valeur maximale ? une valeur minimale ?
\end{enumerate}
\fin
\ex{Fonctions}{b3Jz}{Ex. 7}{5}
Soit $f$ la fonction définie par $f(x)=\dfrac{x\sin x\cos(2x)}{\sqrt{x^2+1}}$.
\begin{enumerate}
\item Montrer que $(\forall x\in\mathbb{R})\ |x|<\sqrt{x^2+1}$.
\item En déduire que la fonction $f$ est bornée.
\end{enumerate}
\fin
\ex{Fonctions}{b3Jz}{Ex. 10}{9}
Soit $f$ la fonction définie sur $\mathbb{R}$ par $f(x)=\dfrac{x^2+1}{x^2+x+1}$.
\begin{enumerate}
\item Montrer que la fonction $f$ est majorée par $2$.
\item Montrer que la fonction $f$ est minorée par $\dfrac23$.
\item Déterminer l'image de $\mathbb{R}$ par la fonction $f$.
\item Montrer que $T_f(x,y)=\dfrac{xy-1}{\left(1+x+x^2\right)\left(1+y+y^2\right)}$.
\item Étudier la monotonie de la fonction $f$ sur chacun des intervalles $]-\infty,-1]$, $[-1,1]$ et $[1,+\infty[$.
\end{enumerate}
\fin
